\section{Summary and Outlook}
One of the central scientific objectives of the R3B experiment is to deliver stringent experimental constraints on the nuclear equation of state (EoS), in particular on the density dependence of the symmetry energy near saturation. A key quantity in this context is the slope parameter $L$, which governs the pressure of neutron-rich matter and thus has direct implications for the neutron-skin thickness of heavy nuclei and for the macroscopic properties of neutron stars. A promising strategy to constrain $L$ is the measurement of neutron-skin thicknesses of exotic, neutron-rich nuclei through neutron-removal cross sections, combined with state-of-the-art reaction models that relate these observables to the underlying symmetry-energy parameters.\\
These reaction models must be benchmarked against systems for which the nuclear density distributions and free nucleon-nucleon cross sections -- the key ingredients of the models -- are known with high accuracy. The total interaction cross section of the $^{12}$C+$^{12}$C system represents an ideal test case, since the charge radius and density distribution of $^{12}$C are well established from decades of scattering experiments.\\
The measurements presented in this thesis of the total interaction cross section for $^{12}$C+$^{12}$C collisions in the energy range of 400–800 MeV/nucleon demonstrate agreement with existing data and theoretical predictions at intermediate energies of 550 and 650 MeV/nucleon. At 400 MeV/nucleon a slight enhancement is observed, while at 800 MeV/nucleon a significant reduction of more than 4\% with respect to theoretical expectations is confirmed, consistent with the findings of Ponnath et al.~\cite{ponnath2024measurement}. The observed increase in transparency at high energies may be related to structural effects such as $\alpha$-clusterization in carbon nuclei, or to in-medium modifications arising from the onset of inelastic channels, including one- and two-pion production, which open at approximately 280 and 530 MeV/nucleon, respectively. Future dedicated measurements will be required to disentangle these contributions and to achieve a consistent description of the reaction mechanism and its energy dependence.
The precision achieved in this work—better than 1\% - demonstrates that the R3B setup is capable of performing measurements with uncertainties below the 2\% threshold required to constrain the symmetry-energy slope parameter $L$ at the $\pm$10 MeV level~\cite{aumann2017peeling}. Building on this result, future experiments with tin and lead isotopes over a wide range of neutron-to-proton ratios are planned to advance the program of constraining the symmetry energy through neutron-removal observables.\\
The second part of this thesis focused on quasi-free scattering (QFS) as an experimental method to probe single-particle structure in nuclei. The feasibility study of the exclusive reaction $^{12}$C(p,2p)$^{11}$B demonstrated that a clean selection of this channel is achievable and that all reaction products can be measured kinematically completely within the R3B setup. The subsequent S455 experiment extended this method to study fission via QFS in the reaction $^{238}$U(p,2p)fission. The analysis conducted in this work shows that QFS-induced fission provides a powerful tool to directly relate fission-fragment yields and distributions to the corresponding excitation energies of the fissioning nucleus. Complementary measurements of Coulomb-excitation-induced fission further revealed the existence of a new “island of asymmetric fission,” driven by shell effects in the light fragments~\cite{morfouace2025asymmetric}. Ongoing and future proposals aim to refine the boundaries of this island and, in the case of QFS-induced fission, to extend these studies to exotic nuclei far from stability.\\
The forthcoming move of the R3B setup to the new High-Energy Cave (HEC) at FAIR will significantly enhance the experimental capabilities. The larger experimental area will enable an optimized detector geometry, while the continued development of detector systems -- together with advanced data-analysis techniques such as machine learning, as demonstrated in this thesis for cluster reconstruction in the CALIFA calorimeter -- will further improve precision and efficiency. These advancements will be essential for the next generation of high-accuracy measurements needed to unveil the structure of exotic nuclei, to refine the nuclear inputs for r-process modeling, and ultimately to improve our understanding of the origin of the heavy elements in the universe.
